\documentclass[review, nonacm, acmsmall]{acmart}

\setcopyright{none}
\citestyle{acmauthoryear}
\usepackage{booktabs}
\usepackage{graphicx}
\usepackage{cleveref}
\usepackage{pgfgantt}

\usepackage[T1]{fontenc}
\usepackage[utf8]{inputenc}
\usepackage{listings}
\usepackage{color}
\crefname{lstlisting}{listing}{listings}
\Crefname{lstlisting}{Listing}{Listings}

\lstset{
  language=Java,
  basicstyle=\ttfamily\small,
  keywordstyle=\bfseries,
  commentstyle=\itshape\color[rgb]{0.4,0.4,0.4},
  columns=fullflexible,
}

\title{Solving Instead of Guessing: Analysis-Derived Inputs for Finding Errors in JPAMB}

% TODO: fill in names, emails and student numbers.
\author{Student One}
\email{sXXXXXX@dtu.dk}
\author{Student Two}
\email{sXXXXXX@dtu.dk}
\author{Student Three}
\email{sXXXXXX@dtu.dk}
\affiliation{%
  \institution{DTU Compute - Group ABE}
  \city{Kgs. Lyngby}
  \country{Denmark}}

\begin{document}

\maketitle

\section{Introduction}

\paragraph{Context}
A program analysis answers questions about a program without the developer
having to run every input by hand: can this method divide by zero, can it
fail an assertion, can it loop forever? The JPAMB benchmark~\citep{jpamb}
asks exactly these questions for a suite of small Java methods, and scores
an analysis by how well it predicts which outcomes can happen.

The two families of analysis in this course answer these questions from
opposite sides. A \emph{static analysis} based on abstract
interpretation~\citep{cousot1977} over-approximates every execution, so when
it never reaches an outcome, the outcome is impossible. When it does reach
the outcome, it cannot tell a real error from an artefact of the
abstraction. A \emph{dynamic analysis} runs the program on concrete inputs,
so every error it observes is real and comes with a witness input, but when
it observes nothing, it has proven nothing.

\begin{lstlisting}[float, label=lst:example, caption={Two JPAMB methods whose errors are guarded by specific values.}]
public static int divideByNMinus10054203(int n) {
  return 1 / (n - 10054203);       // divides by zero only for n = 10054203
}

public static void divisionLoop(int n) {
  for (int i = 0; i < 5; i++) {
    n /= 2;
  }
  assert n != 32;                  // fails only for 1024 <= n <= 1055
}
\end{lstlisting}

\paragraph{Gap}
The errors that matter are usually guarded by specific values, as in
\cref{lst:example}. For the first method, a static analysis with intervals
sees that the divisor may be zero, but can only answer ``maybe''. A random
fuzzer tries values uniformly from the $2^{32}$ possible integers and hits
\texttt{10054203} with probability $2^{-32}$ per run, while testing small
values close to zero never hits it. Neither family produces a confident
``yes''.

The second method shows a further problem with relying on the benchmark
score alone. JPAMB takes its ground truth from the inputs listed in each
case, and for \texttt{divisionLoop} it lists only $n = 0$ and $n = 1$, both
of which terminate normally. The assertion error for $n = 1024$ is real, but
the benchmark counts it as impossible. An analysis can therefore score well
by fitting the listed cases, for example a syntactic analysis that matches
patterns in the source, without its answers being backed by any evidence.
High JPAMB scores are achievable this way, which makes the score a weak
signal on its own.

\paragraph{Innovation}
We propose to stop guessing inputs, and instead \emph{derive} them from the
program, using one analysis to inform the next:
\begin{enumerate}
  \item A \textbf{syntactic analysis} extracts the constants that appear in
    the bytecode, such as \texttt{10054203} and \texttt{32}.
  \item An \textbf{unbounded static analysis} with an interval abstraction
    and widening runs to a fixpoint. Every outcome it never reaches is
    reported as a sound ``no'', and is removed from the search.
  \item For every outcome still open, a \textbf{target-directed input
    generator} locates the instruction that can cause it (a division, an
    array access, or the comparison guarding an assertion), and walks back
    through the arithmetic between a parameter and that instruction,
    inverting addition, subtraction, multiplication and division to solve
    for the input that reaches the boundary. For \cref{lst:example}, this
    gives $n = 10054203$ for the first method and $n = 32 \cdot 2^5 = 1024$
    for the second. Where inversion is not possible, the generator falls back
    to the extracted constants, their neighbours $c-1, c, c+1$, and boundary
    values such as $0$, $\pm 1$ and the ends of the integer range.
  \item A \textbf{concrete interpreter} runs the generated inputs. An input
    that triggers the outcome is a witness, and the answer is a confident
    ``yes''.
\end{enumerate}
The result is deterministic, and every confident answer is backed by either
a proof (``no'') or a witness (``yes'').

We will answer the following research questions:
\begin{description}
  \item[RQ1] Do analysis-derived inputs find error-triggering inputs more
    reliably, and with fewer executions, than random fuzzing and small-value
    testing?
  \item[RQ2] How much of the search does static pruning remove?
  \item[RQ3] Does the approach hold up on programs it was not developed
    against, compared with a syntactic analysis?
\end{description}

Our contributions will be:
\begin{itemize}
  \item an interval abstract interpreter with widening for JPAMB,
  \item a target-directed input generator that inverts arithmetic to solve
    for error-triggering inputs,
  \item a pipeline in which syntactic and static results guide the dynamic
    analysis,
  \item a set of new benchmark methods with errors guarded by specific
    values, and
  \item an empirical comparison against random fuzzing, small-value testing
    and syntactic analysis.
\end{itemize}

\paragraph{Evaluation}
We evaluate on two benchmark sets: the JPAMB suite, and a held-out set of
around 20 new methods that we write \emph{before} implementing the input
generator, with errors behind magic constants, chains of arithmetic, and
loops that transform the input. We will submit some of them to JPAMB.

We compare the full pipeline against five baselines: the syntactic analysis,
random fuzzing, small-value testing (our dynamic analysis), the generator
using only extracted constants, and the full pipeline without static
pruning. Since the JPAMB score can saturate, we also measure the number of
executions before the first witness, the number of answers left as
``maybe'', the share of answers backed by a proof or a witness, and the
running time. Random baselines are repeated 10 times, and we report the mean
and standard deviation. Witnesses that contradict the benchmark's ground
truth, as for \texttt{divisionLoop}, are reported separately.

The main threats to validity are that both benchmark sets are small and
synthetic, and that we write the held-out set ourselves. We reduce the
second by writing the held-out methods before the generator exists, and by
not changing the generator after seeing its results on them.

\bibliographystyle{ACM-Reference-Format}
\bibliography{bibfile}

\appendix

\section{Plan}

\Cref{tab:contributions} lists the technical contributions and how the
points are distributed between us. \Cref{fig:timeline} shows the plan. We
first write the held-out benchmark methods, so they are fixed before the
input generator exists. We then build the interval analysis and the
generator in parallel, join them into the pipeline, and run the evaluation.
We aim to have new benchmark methods accepted into JPAMB before the 16th of
November.

% TODO: check the paper and presentation dates on LEARN.
\begin{figure}[h]
\centering
\caption{Project Timeline}
\label{fig:timeline}
\begin{ganttchart}[
    hgrid,
    time slot format=isodate,
    today=2026-10-19,
    expand chart=\linewidth,
    vgrid={*5{white}, *2{dotted}},
    y unit chart=0.5cm,
    group label font=\bfseries\small,
    bar label font=\footnotesize,
    milestone label font=\bfseries\footnotesize,
    vrule/.style={very thick, red},
    vrule label font=\footnotesize\bfseries
  ]{2026-10-12}{2026-12-14}
  \gantttitlecalendar{month=shortname} \\
  % --- Phase 1: benchmark and analyses ---
  \ganttbar{Write held-out benchmark methods}{2026-10-20}{2026-10-30} \\
  \ganttbar{Submit methods to JPAMB}{2026-10-28}{2026-11-13} \\
  \ganttbar{Constant extraction}{2026-10-20}{2026-10-27} \\
  \ganttbar{Interval analysis with widening}{2026-10-20}{2026-11-06} \\
  \ganttbar{Target-directed input generator}{2026-10-27}{2026-11-13} \\
  % --- Phase 2: integration and evaluation ---
  \ganttbar{Integrate into pipeline}{2026-11-09}{2026-11-20} \\
  \ganttbar{Run evaluation and statistics}{2026-11-16}{2026-11-27} \\
  % --- Phase 3: writing ---
  \ganttbar{Read and relate to papers}{2026-10-26}{2026-11-20} \\
  \ganttbar{Write paper}{2026-11-09}{2026-12-01} \\
  \ganttbar{Record video}{2026-12-02}{2026-12-07}

  \ganttvrule[vrule label node/.append style={anchor=north west}]{Proposal}{2026-10-19}
  \ganttvrule{JPAMB}{2026-11-16}
  \ganttvrule{Paper}{2026-12-01}
  \ganttvrule{Present}{2026-12-10}
\end{ganttchart}
\end{figure}

% TODO: replace S1-S3 with names and student numbers.
\begin{table}[h]
    \centering
    \caption{Contributions, $S_1$ is Student One (sXXXXXX), $S_2$ is Student Two (sXXXXXX), $S_3$ is Student Three (sXXXXXX).}
    \begin{tabular}{@{}llrrrr@{}} \toprule
         ID & Contribution & P & $S_1$ & $S_2$ & $S_3$ \\
         \midrule
         \texttt{ISY} & Implement Syntactic Analysis (constant extraction) & 5 & 5 & 0 & 0 \\
         \texttt{IAB} & Implement Novel Abstractions (intervals) & 5 & 0 & 5 & 0 \\
         \texttt{IBA} & Implement Unbounded Static Analysis & 7 & 0 & 7 & 0 \\
         \texttt{NAN} & Integrate Analyses (3) & 10 & 4 & 2 & 4 \\
         \texttt{EXB} & Extend Benchmark Suite & 5 & 1 & 1 & 3 \\
         \texttt{EST} & Do Statistical Analysis of Results & 5 & 0 & 0 & 5 \\
         \texttt{PWP} & Write Paper & 10 & 4 & 3 & 3 \\
         \texttt{PEX} & Find Illustrating Examples & 4 & 2 & 0 & 2 \\
         \texttt{PDI} & Create a Process Diagram & 3 & 3 & 0 & 0 \\
         \texttt{PRW} & Read and relate to recent papers & 6 & 1 & 2 & 3 \\
         \midrule
         \multicolumn{2}{l}{Total} & 60 & 20 & 20 & 20 \\
         \multicolumn{2}{l}{Suggested Max Grade} & - & 12 & 12 & 12 \\
         \bottomrule
    \end{tabular}
    \label{tab:contributions}
\end{table}

\end{document}
